\chapter{数列极限的一般性质}
\label{chap:sequence-limit-properties}

\begin{chapterguide}
本章从 \(\varepsilon\)-\(N\) 定义建立极限的稳定性：极限唯一，收敛列有界，严格符号可传给充分靠后的项，有限次代数运算与绝对值可以通过极限。证明反复使用误差分配、尾部起点取最大值以及分母远离零。每条性质都同时检查逆命题或严格不等式版本是否成立。
\end{chapterguide}

\section{极限的唯一性}

\begin{theorem}[极限唯一性]
若实数列 \((a_n)\) 同时收敛到 \(a,b\in\R\)，则 \(a=b\)。
\end{theorem}

\begin{proof}
反设 \(a\ne b\)，令 \(\varepsilon=\abs{a-b}/3\)。分别由 \(a_n\to a\) 与 \(a_n\to b\) 取得起点 \(N_1,N_2\)。当 \(n\ge\max\{N_1,N_2\}\) 时，
\[
 \abs{a-b}\le\abs{a-a_n}+\abs{a_n-b}
 <2\varepsilon=\frac23\abs{a-b},
\]
矛盾。
\end{proof}

几何上，半径小于 \(\abs{a-b}/2\) 的两个邻域不相交，而一个完整尾部不能同时包含在两个不交邻域中。唯一性依赖值域距离能够分离不同点；以后在一般拓扑空间中，这对应 Hausdorff 条件。

\begin{corollary}
若数列的两个子列分别收敛到不同极限，则原列发散。
\end{corollary}

\section{收敛数列的有界性及逆命题失败}

\begin{theorem}
每个收敛实数列都有界。
\end{theorem}

\begin{proof}
设 \(a_n\to a\)。对误差 \(1\) 取 \(N\)，则 \(n\ge N\) 时
\(\abs{a_n}\le\abs a+1\)。令
\[
 M=\max\{\abs{a_1},\ldots,\abs{a_{N-1}},\abs a+1\},
\]
便有 \(\abs{a_n}\le M\) 对所有 \(n\) 成立。
\end{proof}

有限个首项必须单独并入界中。极限定义只控制尾部，而“数列有界”要求所有项同时受控。

逆命题不成立。数列 \((-1)^n\) 有界但发散。完整的恢复条件需要额外结构：单调有界列由\chapterref{chap:monotone-recursive-sequences}定理收敛；一般有界列只能由\chapterref{chap:subsequences-bw}保证存在收敛子列。

\begin{proposition}[零列乘有界列]
若 \(a_n\to0\)，而 \((b_n)\) 有界，则 \(a_nb_n\to0\)。
\end{proposition}

\begin{proof}
取 \(M\ge1\) 使 \(\abs{b_n}\le M\)。给定 \(\varepsilon>0\)，最终有
\(\abs{a_n}<\varepsilon/M\)，于是
\(\abs{a_nb_n}<\varepsilon\)。
\end{proof}

有界性不能替换成“每一项都是有限实数”，后者对所有实数列都成立而不给统一控制。例如 \(a_n=1/n\)、\(b_n=n\) 时乘积恒为 \(1\)。

\section{保号性、严格保号与非严格不等式}

\begin{theorem}[局部保号]
若 \(a_n\to a\) 且 \(a>0\)，则存在 \(N\) 使
\[
 a_n>\frac a2>0\qquad(n\ge N).
\]
若 \(a<0\)，则 \(a_n<a/2<0\) 最终成立。
\end{theorem}

\begin{proof}
对 \(a>0\) 取误差 \(\varepsilon=a/2\)，则最终
\(\abs{a_n-a}<a/2\)，从而 \(a_n>a-a/2=a/2\)。负数情形对 \(-a_n\) 应用正数结论。
\end{proof}

极限为零时不能推出尾部固定符号；数列 \((-1)^n/n\to0\) 无限次改变符号。反向命题“若最终 \(a_n>0\)，则极限 \(a>0\)”也不成立，反例为 \(a_n=1/n\)。能稳定传给极限的是非严格不等式。

\begin{theorem}[序关系的闭性]
若 \(a_n\to a\)、\(b_n\to b\)，且最终 \(a_n\le b_n\)，则 \(a\le b\)。特别地，若 \(a_n\ge0\) 最终成立，则 \(a\ge0\)。
\end{theorem}

\begin{proof}
若 \(a>b\)，令 \(\varepsilon=(a-b)/3\)。充分大时
\[
 a_n>a-\varepsilon>b+\varepsilon>b_n,
\]
与最终 \(a_n\le b_n\) 矛盾。
\end{proof}

即使每个 \(n\) 都有 \(a_n<b_n\)，极限也可能相等，例如 \(a_n=0\)、\(b_n=1/n\)。若要得到 \(a<b\)，需要最终存在固定间隙，例如 \(b_n-a_n\ge\delta>0\)。

\section{夹逼定理}

\begin{theorem}[夹逼定理]
设实数列满足最终
\[
 a_n\le b_n\le c_n.
\]
若 \(a_n\to L\)、\(c_n\to L\)，则 \(b_n\to L\)。
\end{theorem}

\begin{proof}
给定 \(\varepsilon>0\)，把两个极限的起点与比较成立的起点取最大值。对充分大的 \(n\)，
\[
 L-\varepsilon<a_n\le b_n\le c_n<L+\varepsilon,
\]
故 \(\abs{b_n-L}<\varepsilon\)。
\end{proof}

\begin{corollary}
若 \(\abs{b_n-L}\le r_n\) 最终成立且 \(r_n\to0\)，则 \(b_n\to L\)。
\end{corollary}

夹逼两侧必须趋于同一极限。若只知 \(a_n\to A\)、\(c_n\to C\) 且 \(A<C\)，最多能推出任何收敛子列极限位于 \([A,C]\)，不能保证 \(b_n\) 收敛。

\section{和、差、积、商的极限定理}

\begin{theorem}[极限的代数运算]
若 \(a_n\to a\)、\(b_n\to b\)，且 \(lambda,\mu\in\R\)，则
\[
 \lambda a_n+\mu b_n\to\lambda a+\mu b,
 \qquad a_nb_n\to ab.
\]
若再有 \(b\ne0\)，则 \(b_n\ne0\) 最终成立，并且
\[
 \frac{a_n}{b_n}\longrightarrow\frac ab.
\]
\end{theorem}

\begin{proof}
线性组合由
\[
 \abs{\lambda(a_n-a)+\mu(b_n-b)}
 \le\abs\lambda\abs{a_n-a}+\abs\mu\abs{b_n-b}
\]
及误差分配得到。

对乘积写成
\[
 a_nb_n-ab=a_n(b_n-b)+b(a_n-a).
\]
收敛列 \((a_n)\) 有界，设 \(\abs{a_n}\le M\)。给定 \(\varepsilon>0\)，分别使
\[
 \abs{b_n-b}<\frac{\varepsilon}{2(M+1)},\qquad
 \abs{a_n-a}<\frac{\varepsilon}{2(\abs b+1)}.
\]
于是乘积误差小于 \(\varepsilon\)。

若 \(b\ne0\)，局部保号定理给出最终
\(\abs{b_n}\ge\abs b/2\)。因此
\[
 \abs{\frac1{b_n}-\frac1b}
 =\frac{\abs{b_n-b}}{\abs{b_n}\abs b}
 \le\frac{2}{\abs b^2}\abs{b_n-b}\longrightarrow0.
\]
商结论由乘积结论应用于 \(a_n\) 与 \(1/b_n\) 得到。
\end{proof}

分母极限非零保证统一下界，也保证商最终有定义。若 \(b=0\)，商可能有有限极限、无穷极限或发散：例如
\[
 \frac{1/n}{1/n}=1,\qquad
 \frac{1}{1/n}=n,\qquad
 \frac{(-1)^n/n}{1/n}=(-1)^n.
\]
因此不存在只凭“分子、分母都趋零”即可决定商极限的法则。

\begin{corollary}
若 \(P\) 是实系数多项式且 \(a_n\to a\)，则 \(P(a_n)\to P(a)\)。若 \(R=P/Q\) 且 \(Q(a)\ne0\)，则 \(R(a_n)\to R(a)\)。
\end{corollary}

\section{绝对值、序关系与极限运算}

\begin{theorem}
若 \(a_n\to a\)，则
\[
 \abs{a_n}\to\abs a.
\]
反之，\(\abs{a_n}\to\abs a\) 一般不能推出 \(a_n\to a\)。
\end{theorem}

\begin{proof}
反三角不等式给出
\[
 \bigl|\abs{a_n}-\abs a\bigr|\le\abs{a_n-a}\longrightarrow0.
\]
反例取 \(a=1\)、\(a_n=(-1)^n\)，则 \(\abs{a_n}=1\) 恒成立而原列发散。
\end{proof}

\begin{proposition}
若 \(a_n\to a\)、\(b_n\to b\)，则
\[
 \max\{a_n,b_n\}\to\max\{a,b\},\qquad
 \min\{a_n,b_n\}\to\min\{a,b\}.
\]
\end{proposition}

\begin{proof}
利用恒等式
\[
 \max\{x,y\}=\frac{x+y+\abs{x-y}}2,\qquad
 \min\{x,y\}=\frac{x+y-\abs{x-y}}2,
\]
再应用已证明的代数运算与绝对值极限。
\end{proof}

有限次运算的稳定性不能未经证明推广到随 \(n\) 增长的运算次数。例如每个固定 \(k\) 都有 \((1+1/n)^k\to1\)，但指数若取 \(n\)，数列 \((1+1/n)^n\) 不趋于 \(1\)。量词中的“固定对象”与“随指标变化的对象”必须区分。

\section*{本章小结}
\addcontentsline{toc}{section}{本章小结}

极限唯一性排除多个候选；收敛性给出尾部控制并连同有限首项产生全局有界性。非零极限能稳定符号和分母下界，非严格序关系在取极限后保持。代数运算证明的共同机制是误差分解、收敛列有界和有限个尾部起点取最大值。严格不等式及绝对值逆向都需要反例辨析。

\section*{习题}
\addcontentsline{toc}{section}{习题}

\noindent\textbf{配置说明：}本章按II级核心章配置，共24道编号题，含41个独立训练单元，建议总用时5至7小时。\exerciserange{9.1}{9.16}为核心必做范围。

\subsection*{A组　唯一性、有界性与符号}
\begin{enumerate}[label=\textbf{9.\arabic*.}]
 \exerciseitem{9.1} \mustdo 用两个不交邻域重新证明极限唯一性。
 \exerciseitem{9.2} \mustdo 证明收敛数列的任意子列收敛到同一极限。
 \exerciseitem{9.3} \mustdo 证明：若 \(a_n\to a\)，则 \(sup_n\abs{a_n}\) 是有限实数；指出确界公理的使用处。
 \exerciseitem{9.4} \mustdo 给出有界发散数列，并证明它不收敛而不是只写出公式。
 \exerciseitem{9.5} \mustdo 若 \(a_n\to a>0\)，证明存在 \(delta>0\) 使 \(a_n\ge\delta\) 最终成立；讨论 \(a=0\) 时的反例。
 \exerciseitem{9.6} \mustdo 若 \(a_n\ge0\) 最终成立且 \(a_n\to a\)，证明 \(a\ge0\)。
 \exerciseitem{9.7} \mustdo 若 \(a_n<b_n\) 对所有 \(n\) 成立，构造极限相等的例子；再给出保证极限严格不等的附加条件。
\end{enumerate}

\subsection*{B组　运算与估计}
\begin{enumerate}[label=\textbf{9.\arabic*.},resume]
 \exerciseitem{9.8} \mustdo 从定义证明 \(a_n\to a\)、\(b_n\to b\) 蕴含 \(a_n-b_n\to a-b\)。
 \exerciseitem{9.9} \mustdo 在乘积极限定理中改用分解
 \[
 a_nb_n-ab=(a_n-a)(b_n-b)+a(b_n-b)+b(a_n-a)
 \]
完成证明，并说明各项怎样控制。
 \exerciseitem{9.10} \mustdo 若 \(a_n\to a\ne0\)，证明 \(1/a_n\) 最终有定义且收敛到 \(1/a\)。
 \exerciseitem{9.11} \mustdo 分别构造 \(a_n,b_n\to0\)，使 \(a_n/b_n\) 收敛到给定 \(L\in\R\)、趋于 \(+\infty\)、以及发散。
 \exerciseitem{9.12} \mustdo 证明零列乘有界列为零列；构造“零列乘无界列”分别趋零、趋非零常数和发散的例子。
 \exerciseitem{9.13} \mustdo 证明若 \(a_n\to a\)，则 \(a_n^k\to a^k\)（固定 \(k\in\N\)）。
 \exerciseitem{9.14} \mustdo 若 \(a_n\to a\ge0\) 且每个 \(a_n\ge0\)，证明 \(\sqrt{a_n}\to\sqrt a\)。分别处理 \(a>0\) 与 \(a=0\)。
 \exerciseitem{9.15} \mustdo 证明 \(\max\{a_n,b_n\}\) 与 \(\min\{a_n,b_n\}\) 的极限定理。
 \exerciseitem{9.16} \mustdo 证明若最终 \(a_n\le b_n\le c_n\) 且两端同趋于 \(L\)，则中间列也趋于 \(L\)。
\end{enumerate}

\subsection*{C组　反例与综合项目}
\begin{enumerate}[label=\textbf{9.\arabic*.},resume]
 \exerciseitem{9.17} \mustdo 判断并证明或反驳：
 \begin{enumerate}[label=(\arabic*)]
  \item \(\abs{a_n}\to0\Rightarrow a_n\to0\)；
  \item \(\abs{a_n}\to1\Rightarrow a_n\) 收敛；
  \item \(a_n^2\to a^2\Rightarrow a_n\to a\)。
 \end{enumerate}
 \exerciseitem{9.18} \mustdo 若 \(a_n\to a\)、\(b_n\to b\) 且 \(a_n\le b_n\) 对无穷多个 \(n\) 成立，能否推出 \(a\le b\)？证明结论或给反例。
 \exerciseitem{9.19} \projectdo\ 【前置：本章全部定理；预计用时：60分钟】证明任何由有限次四则运算、绝对值、最大值和最小值组成的表达式，只要所有分母的极限非零，就可以逐项通过极限。
 \exerciseitem{9.20} \optionaldo 设 \(p_n\) 为偶数时 \(a_n\)，奇数时 \(b_n\)。刻画 \(p_n\) 收敛的充要条件。
 \exerciseitem{9.21} \optionaldo 若 \(a_n\to a\)、\(a\ne0\)，证明 \(operatorname{sgn}(a_n)=\operatorname{sgn}(a)\) 最终成立；说明符号函数在 \(0\) 处为何不能通过极限。
 \exerciseitem{9.22} \optionaldo 设 \(a_n,b_n>0\)、\(a_n/b_n\to1\)。证明 \(a_n-b_n\to0\) 需要额外的有界性条件，并给出无此条件时的反例。
 \exerciseitem{9.23} \opendo 比较“代数运算连续”与本章逐项证明之间的关系。说明为何在连续性尚未定义时，不能用后文定理替代当前证明。
 \exerciseitem{9.24} \projectdo\ 【前置：第7--9章；预计用时：75分钟】建立有限维向量数列的分量判据：\((x_n^{(1)},\ldots,x_n^{(d)})\) 在最大范数下收敛，当且仅当每个分量收敛；给出误差在维数 \(d\) 下的依赖。
\end{enumerate}

\section*{部分习题提示}
\begin{itemize}
 \item \exerciseref{9.3}：先由收敛性证明绝对值集合有上界，再用确界公理取得上确界。
 \item \exerciseref{9.11}：可取 \(a_n=L/n,b_n=1/n\)；其余两种分别调整分子符号或量级。
 \item \exerciseref{9.14}：\(a>0\) 时有理化，\(a=0\) 时用 \(\sqrt{a_n}<\varepsilon\Longleftrightarrow a_n<\varepsilon^2\)。
 \item \exerciseref{9.18}：只需抽取满足比较关系的子列。
 \item \exerciseref{9.22}：取 \(a_n=n+1,b_n=n\)。
\end{itemize}
